\chapter{Particle Localization}\label{ch:localization}

Before the velocity can be interpolated to a particle, the particle must be located on the mesh:
the engine must identify the nearest node and, for most interpolation methods, the element that
contains the position. On a regular latitude--longitude grid this is an indexing operation; on an
unstructured mesh it is a search, performed for every particle at every evaluation site --- the
start of a step, every Runge--Kutta stage (\cref{ch:integration}) and the final position. This
chapter describes the nearest-node search by spatial hashing and the guarantee that the node it
returns is the nearest one, the automatic choice of the hash cell, the containing-element test, and
the alternative of tracking the containing element across edges.

\section{Nearest node by spatial hashing}\label{sec:loc-nearest}

A linear scan over all $N$ nodes would make every evaluation cost $O(N)$. Instead the engine builds a
uniform \emph{grid hash} once, when the mesh is loaded, and scans only the nodes of a small window of
cells around the query (\cref{fig:gridhash}). \index{grid hash}

\paragraph{Construction.} Node longitudes are first placed on one branch --- $[0^\circ,360^\circ)$ or
$[-180^\circ,180^\circ)$, whichever gives the smaller span; a mesh whose span exceeds $350^\circ$ on
both branches is treated as cyclic in longitude. With cell side $c$, the grid starts one cell below
the southernmost node, $\phi_{\min}=\min\phi-c$, and has $n_y=\lceil(\max\phi+c-\phi_{\min})/c\rceil$
rows; a cyclic grid starts at $\lambda_{\min}=0$ with $n_x=\lceil 360/c\rceil$ columns, whose last
(``seam'') cell is only $360-(n_x-1)c$ wide, and a non-cyclic grid carries one padding cell on each
side. A query at $(\lambda,\phi)$ falls in the cell
\begin{equation}\label{eq:loc-cell}
  i=\left\lfloor\frac{(\lambda-\lambda_{\min})\bmod 360}{c}\right\rfloor,\qquad
  j=\left\lfloor\frac{\phi-\phi_{\min}}{c}\right\rfloor,
\end{equation}
clamped to the grid; on a non-cyclic grid a query outside the domain that lies angularly closer to the
western edge is folded to the negative side first, so the clamp picks the nearer edge. The nodes of every cell are stored contiguously (a compressed offset array and
a node-index array), so a row of cells is one contiguous slice.

\paragraph{Distance.} Distances are compared as squared chord lengths between points on the unit
sphere, $\lVert\mathbf{r}_p-\mathbf{r}_n\rVert^2$ with
$\mathbf{r}=(\cos\phi\cos\lambda,\;\cos\phi\sin\lambda,\;\sin\phi)$. Minimising the chord minimises the
great-circle distance, without the dateline discontinuity or the polar distortion of a planar
metric. The node vectors are computed once, when the hash is built, into a table of three numbers
per node, so a scan reads coordinates instead of evaluating trigonometric functions per candidate;
the table is filled with exactly the arithmetic the scans used before, so the nearest node found is
bit-for-bit unchanged, ties included (the first node found wins a tie) --- provided Numba runs without
Intel's vector math library, which the hash build checks on sampled nodes and warns about. The plain three-by-three
window scan returns the nearest node among the nodes of the query cell and its eight neighbours.

\begin{figure}[t]
\centering
\begin{tikzpicture}[scale=0.95]
  \fill[lagblue!10] (1,1) rectangle (4,4);
  \draw[help lines,gray!50] (0,0) grid (6,5);
  \foreach \x/\y in {0.4/0.6,1.3/2.4,2.6/3.3,3.4/1.7,2.2/1.2,4.7/2.1,5.3/3.6,0.8/4.2,3.8/4.4,5.1/0.7,1.7/3.7,4.3/3.2}
     \node[meshnode] at (\x,\y) {};
  \node[drifter] (D) at (2.7,2.4) {};
  \node[meshnode,label={[nodelbl]above:nearest}] (NN) at (2.6,3.3) {};
  \draw[velarrow,densely dashed] (D)--(NN);
  \node[font=\footnotesize,lagblue!60!black,anchor=south west] at (1.0,4.02) {3$\times$3 search window};
  \node[drifter,label={[font=\footnotesize]right:query position}] at (7.0,3.2) {};
  \node[meshnode,label={[font=\footnotesize]right:mesh node}] at (7.0,2.4) {};
\end{tikzpicture}
\caption{Nearest-node search by spatial hashing. Nodes are bucketed into uniform cells when the mesh
is loaded. A query (the drifter) is first matched against the nodes of its own cell and the eight
neighbouring cells (shaded); \cref{sec:loc-exact} gives the test that decides whether this window is
enough.}
\label{fig:gridhash}
\end{figure}

\codeanchor{Construction \fn{build\_grid\_hash} (with \fn{\_gh\_lon\_branch},
\fn{\_gh\_grid\_geometry}, \fn{\_gh\_cell\_of}); node table \fn{\_node\_xyz\_fill} and its self-check
\fn{\_node\_xyz\_selfcheck}; window scan \fn{\_grid\_hash\_nearest\_3x3}; batch query
\fn{find\_nearest\_nodes}; all in \srcfile{core/mesh.py}.}

\section{Choosing the cell}\label{sec:loc-cell}

The reader of each mesh family proposes a cell of twice the mean node spacing,
\begin{equation}\label{eq:loc-legacy}
  c_{\text{legacy}}=\max\!\left(\ang{0.05},\;2\sqrt{A/N}\right),
\end{equation}
with $A$ the area of the mesh's longitude--latitude bounding box in square degrees. On a
quasi-uniform mesh this puts a few dozen nodes in the window at every resolution. On a mesh whose
resolution varies by orders of magnitude it does not: a cell sized for the coarse part holds
thousands of nodes where the mesh is fine, and the scan turns into a linear pass over them.

Every run therefore chooses the cell by measurement: the key
\conf{mesh.grid\_hash\_cell\_rule} defaults to \conf{occupancy}. The rule samples up to \num{200000} element centres, finds each one's exact
nearest node once, and evaluates the candidate cells $c_k=c_{\text{legacy}}/\sqrt{2}^{\,k}$,
$k=0,\dots,12$, without rebuilding the hash: for every sample it counts the nodes the window holds,
tests whether the window proves its answer (\cref{sec:loc-exact}), and estimates the cost of the query
--- the nodes scanned, the cells visited and one certificate per widening, a certificate weighing as
much as twenty cells. With the exact-anchor guard of \cref{sec:loc-exact} on (the default), the rule keeps
the reader's cell unless the cheapest eligible candidate (at most 20\,\% of the queries widened, at most
\SI{128}{MB} of offsets) is predicted to be at least 2.75 times cheaper; with the guard off it takes the
smallest cell that leaves no window empty, holds a median of at least eight nodes per window and has at
most half a percentage point more approximate anchors than the reader's cell. The ladder stops early where a candidate would need more than 64 rows of widening
budget or too many cells. With the guard on, quasi-uniform meshes keep the reader's cell and structure-resolving
meshes take a smaller one
(\cref{sec:perf-anchor}). An explicit \conf{mesh.grid\_hash\_cell\_size} overrides the rule, and
\conf{grid\_hash\_cell\_rule: legacy} keeps the reader's cell. The resolved cell and the rule are
written to the output metadata, so any run can be repeated with its cell pinned.

\codeanchor{\fn{choose\_grid\_hash\_cell} (the per-sample statistics in \fn{\_gh\_sample\_stats}),
\fn{apply\_grid\_hash\_cell\_rule}, \fn{set\_grid\_hash\_cell} and the one-line report
\fn{describe\_grid\_hash}, in \srcfile{core/mesh.py}; applied by \fn{run\_tracking} after the mesh is
loaded.}

\section{Exact nearest nodes}\label{sec:loc-exact}

A non-empty three-by-three window does not guarantee that the node it returns is the nearest one: a
closer node can lie just outside the window, most easily at high latitude, where a degree of longitude
shrinks, or with a cell smaller than the local node spacing. A near-nearest anchor changes the
interpolation stencil and, over time, the trajectory. With \conf{mesh.exact\_nearest: true} (the
default) the engine proves every anchor.

\paragraph{The certificate.} Let the window return a node at angular distance
$\theta=\arccos(\sin\phi_p\sin\phi_n+\cos\phi_p\cos\phi_n\cos\Delta\lambda)$ from the query. Any closer
node lies inside the spherical cap of radius $\theta$ around the query, whose half-extent is $\theta$
in latitude and
\begin{equation}\label{eq:loc-cap}
  \Delta\lambda_{\text{cap}}=\arcsin\!\left(\frac{\sin\theta}{\cos\phi_p}\right)
\end{equation}
in longitude, unbounded when $|\phi_p|+\theta\ge 90^\circ$ (the cap reaches a pole). The answer is
proven when the cap fits inside the latitude and longitude intervals the window actually searched ---
taking into account the narrower seam cell of a cyclic grid, a window clipped at the grid edge (there
are no nodes beyond it), and, on a non-cyclic grid, the $\pm360^\circ$ images of the cap
(\cref{fig:loc-certificate}). The test returns which sides fail: the rows, the columns, or --- when the
cap reaches a pole or an image of it meets the node range --- all columns.

\paragraph{Widening.} An unproven window grows in the failing direction --- rows and columns
independently, all columns only where required --- and is scanned again, new rows as contiguous slices,
until the best node found is proven. An empty window grows as a square until the first node is found,
or by whole rows in the polar cap, where a row's length $360^\circ\cos\phi$ is at most 32 cell widths. The growth has a budget: a query
whose nearest node is farther away than the longest element edge (plus two cells) lies outside every
element, so it keeps the best node seen and dies in that step; the budget, between 4 and 4096 rows, is
derived from the mesh when the hash is built.

\paragraph{One search for every path.} The batch query and every fused kernel's evaluation site call
the same exact search, so fused and unfused paths return the same anchor bit for bit, ties included,
and the fused kernels need no round trip through the batch path. With the guard off
(\conf{exact\_nearest: false}; together with \conf{grid\_hash\_cell\_rule: legacy} this reproduces the
nearest-node anchors of runs made before September 2026) the batch path keeps the plain window and
repairs empty windows through a $k$-d tree, and the fused kernels widen an empty window ring by ring ---
an approximate anchor. The two-dimensional barycentric kernels, and the least-squares kernels on HEALPix
pixel meshes, are then enabled only on meshes that pass a density certificate (every cell of every
element's bounding box has a non-empty window); the three-dimensional RK2 kernel instead flags a particle
whose window is empty and re-runs it with the anchor of the batch search (\cref{sec:fused-identity}).

\begin{figure}[t]
\centering
\begin{tikzpicture}[scale=0.85]
  \fill[lagblue!10] (1,1) rectangle (4,4);
  \fill[lagelem!10] (4,1) rectangle (5,4);
  \draw[help lines,gray!50] (0,0) grid (7,5);
  \draw[lagblue,line width=0.9pt] (1,1) rectangle (4,4);
  \draw[lagelem,line width=0.9pt,densely dashed] (1,1) rectangle (5,4);
  \node[drifter] (Q) at (3.55,2.5) {};
  \node[meshnode] (F) at (2.35,3.35) {};
  \draw[kelp!70,line width=0.7pt] (Q) ellipse (1.46 and 1.46);
  \draw[densely dotted] (Q)--(F);
  \node[meshnode] (T) at (4.45,1.95) {};
  \node[font=\footnotesize,anchor=west] at (7.2,4.2) {\tikz\draw[lagblue,line width=0.9pt](0,0)rectangle(0.3,0.2); searched 3$\times$3 window};
  \node[font=\footnotesize,anchor=west] at (7.2,3.5) {\tikz\draw[kelp!70,line width=0.7pt](0.15,0.1)circle(0.12); cap of radius $\theta$};
  \node[font=\footnotesize,anchor=west] at (7.2,2.8) {\tikz\draw[lagelem,line width=0.9pt,densely dashed](0,0)rectangle(0.3,0.2); widened by one column};
  \node[font=\footnotesize,anchor=west] at (7.2,2.1) {\tikz\node[meshnode]{}; found node / true nearest};
  \node[font=\footnotesize,fill=white,inner sep=1pt,above left=0pt and 0pt of F] {found};
  \node[font=\footnotesize,fill=white,inner sep=1pt,right=2pt of T] {nearest};
\end{tikzpicture}
\caption{The exactness certificate. The window returns a node at angular distance $\theta$; a closer
node can only lie in the cap of radius $\theta$ around the query. Here the cap crosses the window's
eastern edge, so the window cannot prove its answer; it widens by one column, finds the true nearest
node, and the smaller cap of the new answer fits inside the searched interval.}
\label{fig:loc-certificate}
\end{figure}

\codeanchor{Cap extents \fn{\_cap\_extents}; certificate \fn{\_anchor\_provable\_rect} (and
\fn{\_anchor\_provable} for the plain window); exact search \fn{\_grid\_hash\_nearest\_exact}; budget
\fn{\_gh\_far\_rows\_for}; guard flag \fn{set\_exact\_nearest}; guard-off single-particle search
\fn{\_grid\_hash\_nearest\_widened}; density certificate \fn{grid\_hash\_dense\_for\_elements}; all in
\srcfile{core/mesh.py}. The fused kernels anchor each site through \fn{\_fused\_site\_anchor} in
\srcfile{numerics/integrators.py}.}

\section{Where localization runs}\label{sec:loc-where}

Every evaluation site of the step re-queries the hash; no step path updates the anchor from the previous
one (only the offsets of the diffusivity gradient, \cref{ch:diffusion}, start a local walk from the
particle's node). A fused
kernel anchors its start, its stages and its final position through the exact search and tests the
final position for containment; the unfused paths query the batch search for the particles that
moved and test their containment. A particle folded back into a periodic window
(\cref{ch:lifecycle}) is re-anchored after the fold. On a HEALPix mesh in the geographic frame, points
poleward of the first ring of pixel centres lie outside every element and are lost; the rotated frame
of \cref{ch:frame} removes this limit.

\section{The containing element}\label{sec:loc-contain}

The nearest node does not by itself determine every interpolation stencil; several methods need the
element that \emph{contains} the particle. The engine scans the elements of the nearest node's
one-ring and tests each for containment (\cref{fig:contain}). If none contains the particle --- which
can occur on an irregular mesh, where the nearest node is not always a vertex of the containing
element --- the search widens to a second ring, scanning the elements adjacent to every vertex of the
one-ring elements. A particle for which neither ring yields a containing element is taken to lie
outside the domain; its step is undone and its status set to \code{OUTSIDE} (\cref{ch:lifecycle}).

\begin{figure}[t]
\centering
\begin{tikzpicture}[scale=1.25]
  \coordinate (C)  at (2.00,1.55);
  \coordinate (p1) at (0.35,1.50);
  \coordinate (p2) at (1.15,0.20);
  \coordinate (p3) at (2.95,0.35);
  \coordinate (p4) at (3.70,1.70);
  \coordinate (p5) at (2.75,2.95);
  \coordinate (p6) at (1.00,2.80);
  \foreach \a/\b in {p1/p2,p2/p3,p3/p4,p4/p5,p5/p6,p6/p1}
     \fill[meshface] (C)--(\a)--(\b)--cycle;
  \fill[hilface] (C)--(p3)--(p4)--cycle;
  \foreach \a/\b in {p1/p2,p2/p3,p3/p4,p4/p5,p5/p6,p6/p1}{
     \draw[meshedge] (C)--(\a); \draw[meshedge] (\a)--(\b);}
  \node[drifter] (D) at (2.78,1.35) {};
  \draw[densely dotted,kelp!70] (D)--(C);
  \draw[densely dotted,kelp!70] (D)--(p3);
  \draw[densely dotted,kelp!70] (D)--(p4);
  \foreach \p in {C,p1,p2,p3,p4,p5,p6} \node[meshnode] at (\p) {};
  \node[nodelbl,above left=0pt of C] {nearest node};
  \node[font=\footnotesize,kelp!70!black,below right=-1pt of D] {drifter};
  \node[font=\footnotesize] at (2.95,0.95) {sub-triangles};
\end{tikzpicture}
\caption{Containing-element search. The engine scans the one-ring of the nearest node and tests each
element for containment. The drifter lies in the shaded triangle; the dotted segments split it into
three sub-triangles whose signed areas, normalized by the whole, are the barycentric coordinates that
decide containment.}
\label{fig:contain}
\end{figure}

The containment test proceeds in three steps. First, the valid vertices of the candidate element are
compacted, discarding the $-1$ padding and, on the mixed meshes of FESOM-C, the repeated vertex.
Second, a dateline-safe bounding-box prefilter, in longitude differences wrapped to
$[-180^{\circ},180^{\circ})$, rejects candidates that cannot contain the point. Third, the element and
the point are projected to a local plane --- an equirectangular plane around the mean vertex,
$x=\Delta\lambda\,R\cos\phi_0$, $y=\Delta\phi\,R$, with the mean longitude taken on the wrapped
differences when the element straddles the dateline --- and the point is declared inside when its
barycentric coordinates are all non-negative to within a small negative tolerance, so that a point
exactly on an edge is inside and no particle falls between two neighbours. A near-degenerate triangle,
whose area falls below a second tolerance, is rejected rather than divided by. A quadrilateral is
tested on a fan of triangles, which is exact because quadrilaterals are required to be convex
(checked when the mesh is loaded, \cref{ch:mesh}); its interpolation weights are the Wachspress
coordinates of \cref{sec:hvel-quad}, computed separately.

\codeanchor{The ring search is \fn{\_find\_containing\_element\_local}; the containment test is
\fn{\_element\_contains} (bounding box) and \fn{\_polygon\_contains\_bool} (barycentric signs) on
\fn{\_project\_to\_local\_xy}; the tolerances \code{\_BARY\_EPS} $=-10^{-6}$ and \code{\_AREA\_TOL}
$=10^{-10}$ are module constants. All in \srcfile{core/mesh.py}.}

\section{Tracking the element across edges}\label{sec:loc-walk}

With \conf{numerics.elem\_search: walk} the containing element is searched once, when the run starts
(or taken from the checkpoint of a resumed run), and afterwards carried from step to step: at every evaluation site the engine starts from
the carried element and walks across edges towards the position. In an element with signed area $S$
and compacted vertices $(x_k,y_k)$, the edge coordinates of the point are
\begin{equation}\label{eq:loc-walk}
  w_k=\frac{\operatorname{sgn}S}{|S|}\,\bigl[(x_{k+1}-x_k)(p_y-y_k)-(p_x-x_k)(y_{k+1}-y_k)\bigr],
\end{equation}
and the point is inside when $\min_k w_k\ge-10^{-6}$; otherwise the walk crosses the edge with the
most negative $w_k$ into the neighbouring element. An acceptance within $10^{-2}$ of an edge is
re-checked with the exact containment predicate of \cref{sec:loc-contain}. Every outcome that is not a
clear containment --- a boundary edge, a two-element cycle, a degenerate element, or 256 steps --- falls
back to the ring search of \cref{sec:loc-contain}, so the walk returns the same element as the ring
search wherever that search is exhaustive and finds elements beyond the second ring where it is not.
The carried element is written to checkpoints. On a 28\,000-element mesh the walk costs
\SI{35}{ns} per particle against \SI{126}{ns} for the ring search; it runs on the unfused paths only
(the fused kernels keep their own search), which is why it pays off for element-layout barycentric
and edge least-squares interpolation and is left at \conf{local} on node-layout runs whose fused
kernels would otherwise be switched off.

\codeanchor{\fn{\_walk\_containing\_element} and \fn{find\_containing\_elements\_walk} in
\srcfile{core/mesh.py}; the element neighbours come from \fn{build\_elem\_neighbors}; the tracking
layer seeds and updates the carried elements (\fn{\_walk\_update\_elems} in
\srcfile{numerics/integrators.py}).}
